\section{Prediction under intervention}
\label{sec:C06:intervention}
A world model for an agent predicts consequences conditional on actions and the information available before those actions. Let $e_t$ be environment state, $a_t$ an action, and $o_t$ an observation. A fully observed model estimates $P(e_{t+1}\mid e_t,a_t)$ and a reward or task-progress function. Under partial observability it instead maintains a belief or learned latent state from $H_t$. Predicting a plausible next screenshot is not sufficient: the prediction must distinguish the consequences of the agent's available interventions, including failure, denial, and partial completion. A model that generates visually convincing but causally incorrect transitions can actively mislead a planner.

Observational data alone can leave the intervention effect unidentified. Suppose a training corpus shows a confirmation dialog followed by a completed purchase. The correlation does not establish which button was clicked, whether authorization was granted, or whether an unseen service completed the transaction. A controller using this correlation as an action model can infer success without performing the required operation. Action-labeled transitions, controlled interventions, and authoritative postconditions reduce this ambiguity. They do not eliminate every hidden confound, especially when external services contain state not exposed to the agent.

The model's boundary should include failure semantics. A network timeout is not equivalent to a failed transaction: the remote operation may have committed while the response was lost. A useful transition model can represent this uncertainty and the value of a status query. A text generator trained only on successful demonstrations may instead fill in the expected success narrative. The distinction is operationally important for stateful agents because imagined certainty can cause duplicate execution, incorrect memory updates, or premature termination. The world model should support epistemic uncertainty about the actual external state, not merely stochastic variation in its rendered description.

World modeling also differs from a policy prior. A language model may know that tests usually follow a code edit, which helps propose an action sequence. That does not mean it predicts whether the particular patch passes the repository's tests. A learned policy can exploit regularities without an explicit simulator; a simulator can be useful even when the final policy executes without online lookahead. The comparative question is whether adding the model improves decisions under matched data and computation, and whether the gain comes from predictive accuracy, additional supervision, or extra search.

